% BHU PYQ -- Manually transcribed & error-corrected
\documentclass[11pt,a4paper]{article}

\usepackage[a4paper,top=2.5cm,bottom=2.5cm,left=2.8cm,right=2.8cm,headheight=14pt]{geometry}
\usepackage{amsmath,amssymb}
\usepackage[T1]{fontenc}
\usepackage[utf8]{inputenc}
\usepackage{lmodern}
\usepackage{microtype}
\usepackage{fancyhdr}
\usepackage{enumitem}
\usepackage{hyperref}
\usepackage{lastpage}
\usepackage{parskip}

\hypersetup{colorlinks=true,urlcolor=black,linkcolor=black,
  pdftitle={CHB-604: Physical Chemistry-IV -- BHU PYQ},pdfauthor={Banaras Hindu University}}

\pagestyle{fancy}\fancyhf{}
\renewcommand{\headrulewidth}{0.4pt}\renewcommand{\footrulewidth}{0.4pt}
\fancyhead[L]{\small\textit{Banaras Hindu University}}
\fancyhead[R]{\small\textit{CHB-604: Physical Chemistry-IV}}
\fancyfoot[C]{\small Page \thepage\ of \pageref{LastPage}}

\newlist{parts}{enumerate}{2}
\setlist[parts,1]{label=(\alph*),leftmargin=*,itemsep=5pt,topsep=3pt}
\setlist[parts,2]{label=(\roman*),leftmargin=*,itemsep=3pt,topsep=2pt}
\newcommand{\pts}[1]{\hfill\textit{[#1]}}

\begin{document}

\begin{center}
  {\large\bfseries BANARAS HINDU UNIVERSITY}\\[2pt]
  {\normalsize B.Sc. (Hons.) Semester VI Examination, 2022-23}\\[6pt]
  \rule{0.8\linewidth}{0.5pt}\\[6pt]
  {\large\bfseries Paper No. CHB-604: Physical Chemistry-IV}\\[4pt]
  {\normalsize Time: 3 Hours\hfill Maximum Marks: 70}
\end{center}
\rule{\linewidth}{0.4pt}
\small
\noindent\textit{Instructions: Answer five questions in total, including Question~1 which is compulsory.}
\normalsize
\bigskip

\noindent\textbf{Question 1.} Answer all parts of the following: \pts{5 $\times$ 2 = 10}
\begin{parts}
  \item Characterize the following molecules as linear, spherical tops, symmetrical tops, and asymmetrical tops: $\text{H}_2\text{O}, \text{CH}_4, \text{CH}_3\text{Cl}, \text{COS}, \text{C}_6\text{H}_6$.
  \item Consider the three ensembles: microcanonical, canonical, and grand canonical. Which physical properties are constant in each of them?
  \item What is zero-point energy? Explain.
  \item State the physical significance of the partition function.
  \item Differentiate between chemical shift and coupling constant in NMR.
\end{parts}

\medskip\rule{\linewidth}{0.2pt}

\noindent\textbf{Question 2.}
\begin{parts}
  \item Construct the Hamiltonian operator for a $1$-D harmonic oscillator and set up the Schr\"odinger wave equation for this system. \pts{6}
  \item Discuss the Born-Oppenheimer approximation and its significance in molecular spectroscopy. \pts{8}
\end{parts}

\medskip\rule{\linewidth}{0.2pt}

\noindent\textbf{Question 3.}
\begin{parts}
  \item Schematically represent the different vibrational modes of $\text{CO}_2$ and $\text{H}_2\text{O}$ molecules. Identify which modes are IR-active. \pts{6}
  \item Show that the lines in the rotational spectrum of a diatomic molecule are equally spaced by $2B$. \pts{8}
\end{parts}

\medskip\rule{\linewidth}{0.2pt}

\noindent\textbf{Question 4.}
\begin{parts}
  \item What is Larmor precession? Show that the Larmor precessional frequency is equal to the transition frequency between two nuclear spin states. \pts{6}
  \item State the Franck-Condon principle and explain its potential energy curve diagrams. \pts{8}
\end{parts}

\medskip\rule{\linewidth}{0.2pt}

\noindent\textbf{Question 5.} Answer any two of the following: \pts{2 $\times$ 7 = 14}
\begin{parts}
  \item Postulates of Quantum Mechanics.
  \item Derive an expression for the translational partition function of a molecule of mass $m$ free to move in a $3$-D box of volume $V$.
  \item Ensembles and their thermodynamic relationships.
\end{parts}

\vfill
\rule{\linewidth}{0.4pt}
\begin{center}\small
  \href{https://bhu-exam.vercel.app/}{Click here to visit our website}\\[4pt]
  \href{https://me-aryan.vercel.app/}{Click here to visit our contributor}\\[4pt]
  \href{https://quashan-me.vercel.app/}{Click here to visit our contributor}
\end{center}

\end{document}
